\documentclass[10pt]{article}
\usepackage[a4paper,margin=24mm]{geometry}
\usepackage{fontspec}
\setmainfont{Libertinus Serif}
\usepackage{mathtools}
\usepackage{unicode-math}
\setmathfont{Libertinus Math}
\usepackage[english]{babel}
\usepackage{microtype}
\usepackage[hidelinks,unicode]{hyperref}
\hypersetup{pdftitle={D026 restrained textual and translation apparatus},pdfauthor={Pierre Deligne}}
\providecommand{\dashrightarrow}{\mathrel{⇢}}
\setlength{\parindent}{0pt}
\setlength{\parskip}{0.55em}
\begin{document}
\begin{center}
{\Large\bfseries D026: restrained textual and translation apparatus}\par
\smallskip
Authority pages 1--18; printed folios 299--316
\end{center}
This apparatus records only material readings, logged repairs, translation choices, and semantic diagram decisions. It does not promote inherited comparison work or copy matter into the canonical editions.


\subsection*{Authority page 1 / printed folio 299}

Canonical text and critical pixels. The title, byline, dedication, Contents, and opening formulas were replayed from the controlling raster; the journal masthead, copyright line, folio, and scanner margins were excluded. The dedication retains the source’s unaccented initial “A” and terminal period. Contents item 2 retains “représentations virtuelles”; the actual §2 heading independently reads “représentations réelles”. The compositorial singular “définis en terme des Frobenius géométriques” is corrected to the idiomatic “définis en termes des Frobenius géométriques”. Equation (1.1.3) remains \(\varepsilon(V,s)=W(V)(f\cdot D^{\dim(V)})^{1/2-s}\); the inherited lowercase \(w\) and outer exponent \(s-1/2\) are rejected. No crop.\par

\subsection*{Authority page 2 / printed folio 300}

Canonical text and notation. The additive character is \(\psi\), not the inherited \(\theta\). The defective phrase “nous n’en avons qu’à la représentation” is repaired to “nous n’avons affaire qu’à la représentation”; the English already conveyed that intended sense. The authority’s local reciprocity identification is \(K^*\simeq W(\bar K/K)^{\mathrm{ab}}\); the earlier state’s equality sign is corrected. The global normalization is \(\int_{A_K/K}\bigotimes_v dx_v=1\), and (1.2.1) has left side \(\varepsilon(V,s)\). Equation (1.2.4) is independently \(W(V)f^{1/2-s}\). No crop.\par

\subsection*{Authority page 3 / printed folio 301}

Textual decisions. The page boundary reads \(W(V_v)\), not the inherited \(W(V_0)\). Stiefel-Whitney classes use superscripts \(w^1,w^2\); the codomain in the composite defining \(w^*\) is the unit group marked \({}^\times\). In (1.4.1) the arrows are labelled “infl”, “(1)”, and “inv”. Translation note: “réalisable sur \(\mathbf R\)” is consistently rendered “realizable over \(\mathbf R\)”; \(\operatorname{cl}\) and \(\operatorname{inv}\) remain notation rather than being expanded in prose. No crop.\par

\subsection*{Authority page 4 / printed folio 302}

Textual decisions. The actual section heading is “Lemmes sur les représentations réelles des groupes finis”, distinct from Contents item 2 on page 1. The two commuting squares were reconstructed semantically with all objects and vertical labels (Ind, tr, and \(f_*\)); the inherited candidate’s conflated diagram was rejected. The trace map is \(\operatorname{Tr}_f:f_*\mathbf Z/2_X\to\mathbf Z/2_Y\). Translation note: “revêtement fini” is rendered “finite covering”. No authority crop was needed.\par

\subsection*{Authority page 5 / printed folio 303}

Diagram decisions. The exact-sequence square (2.2.1), the dashed lifting triangle \(\mathfrak S_k\dashrightarrow\operatorname{Spin}(nk)\), and the additive obstruction square were replayed semantically, preserving all arrows, \(\Sigma\), and the vertical map \(a+b\). The local system is \(((\mathbf Z/2)^k)^Q\), and the bundle is \(f_*\mathcal V\). Translation note: “groupe structural” is rendered “structure group”. No crop was required.\par

\subsection*{Authority page 6 / printed folio 304}

Canonical grammar and formula control. The printed English intrusion “Lemma.” is corrected to French “Lemme.”, and “Une lemme analogue” to “Un lemme analogue”. The determinant identity retains \(\det(V+W)\). The authority formula is \(r(\chi)=\operatorname{Ind}_{\mathbf Z/n}^{D_n}([\chi]-[1])\); the inherited \(\operatorname{Ind}(\chi)-[\chi]-[1]\) is rejected. The terminal induction identity is transcribed without an editorial restriction symbol absent from the print. No crop.\par

\subsection*{Authority page 7 / printed folio 305}

Formula decisions. The three-line expansion of \(\sum[\varepsilon_i]\) was reconstructed directly from the authority: its second line has two nested-product tensor terms, and the resulting expression for \(V\) carries the printed \(\pm\). The inherited branch instead introduced pairwise and square-character terms and omitted \(\pm\); it was rejected. The page contains no nonsemantic figure, so no crop was required.\par

\subsection*{Authority page 8 / printed folio 306}

Canonical grammatical repair. The scan reads ‘et seuls \(\chi\) et \(\bar\chi\) figurant dans \(V|_B\)’, which leaves the coordinated subject without a finite verb. The canon adopts ‘figurent’; the English accordingly reads ‘occur’. Induction, restriction, quotient notation, and the page boundary were otherwise preserved. The inherited branch agrees on ‘figurent’ but remains ZERO\_ACCEPTED and was not promoted. No crop.\par

\subsection*{Authority page 9 / printed folio 307}

Canonical repairs and comparison. The scan prints \(\chi\in H^2(\operatorname{Gal}(\bar K/K),\mathbf C^*)\), but the connecting morphism from the exponential sequence necessarily has source \(H^1(\mathbf C^*)\); the canon therefore adopts \(H^1\), with the scan reading preserved in the source record’s repair metadata. The inherited branch happens to agree on this degree but corrupts the preceding conjugation identity, changes \(c^1(W)\), drops the completion hat on \(K^*\), and omits \(\chi\) after \(\operatorname{Ind}\). No crop.\par

\subsection*{Authority page 10 / printed folio 308}

Canonical repairs and formula control. Because \(\chi\) is a character of \(L^*\), the Tate factor on this page must use the additive character \(\psi_L=\psi\circ\operatorname{Tr}_{L/K}\) of \(L\); the scan’s unsubscripted \(\psi\) is therefore restored to \(\psi_L\) in the epsilon-factor lines and in the dependence of \(c\). In the long equality after \(c_1=c\lVert\Delta\rVert^{-1/2}\), the scan prints \(\chi^{-1}(a)\) in an integral whose bound variable is \(x\); the canon adopts \(\chi^{-1}(x)\). Every occurrence of \(f\), \(\widehat f\), \(da\), \(d^*a\), \(a\bar x\), and \(ax\Delta\) was checked independently. The inherited derivation corrupts several of these terms and was rejected. No crop.\par

\subsection*{Authority page 11 / printed folio 309}

Canonical repair and notation. The scan’s undefined ‘triviale sur \(F^*\)’ is corrected to \(K^*\), the field used throughout the quadratic extension \(L/K\) and the quotient \(L^*/K^*\). The obstruction is \(\alpha w_1^2+w_2\), not the inherited \(\alpha w_1+w_2\); the total class ends in \(1+w^2+\cdots\), not a barred class. The Galois spectral-sequence row uses \(\widehat{\mathbf Z}\), while the Weil row uses \(\mathbf Z\). No crop.\par

\subsection*{Authority page 12 / printed folio 310}

Canonical repairs, array replay, and one S03 supersession. Lemma 4.2.4 is printed with the coefficient omitted from the domain; the canon restores \(H^2(W(\bar K/K),\mathbf Z(1))\), required by the extension and spectral-sequence context. The scan says \(E_2^{pq}=0\) for even \(q\), whereas the displayed array has its odd-\(q\) rows zero; the canon adopts “odd”, which is the condition implying \(E_2=E_3\). S03 withdraws the earlier change of \(z\in\bar K\) to \(z\in\bar K^*\): after restriction to (4.2.3), the preimage of \(\bar K^*\) in \(E\) is the additive covering group \(\bar K\), and the lifted conjugation relation is correctly stated for \(z\in\bar K\). This is precisely what makes \(\widetilde F^2\) fail to commute with \(\widetilde F\), contradicting the fact that an element commutes with its square. The four-row \(E_2\) array and \(d_3:E_3^{0,2}\to E_3^{3,0}\) are retained. No crop.\par

\subsection*{Authority page 13 / printed folio 311}

Canonical repairs and typing. The continuation of the real-Weil-group computation gives \(H^2=\mathbf Z\), index two under restriction, and \(H^3=0\). The scan’s singular “des formule” is normalized to “des formules”. More materially, the final displayed reciprocity identity in the scan equates the \(\mathbf Q/\mathbf Z\)-valued invariant directly with the \(\mathbf C^*\)-valued character value. The canon restores the omitted \(\exp(2\pi i\,\cdot)\), as already required by Proposition 4.6 and by the coefficient types. The inherited TeX retained the ill-typed equality and was rejected. No crop was required.\par

\subsection*{Authority page 14 / printed folio 312}

Canonical repairs and local-global notation. In the first paragraph the scan omits \(W\) from the local cohomology group; the canon restores \(H^2(W(\bar K_v/K_v),\mu_n)\), the target of restriction from global Weil cohomology and the group that factors through \(H^2(\mathbf Z,\mu_n)\) at almost all places. The proof of 4.9.1 cites 4.5 in print; the operative exponential invariant formula is Proposition 4.6, so the cross-reference is corrected to 4.6. The barred antecedent \(\bar c\), the modulus notation, \(L_{\mathbf R}^*\), \(\phi\), and \(\mathfrak m'\ge\mathfrak m\) follow the authority pixels. No crop.\par

\subsection*{Authority page 15 / printed folio 313}

Diagram replay. The definition of \((\mathbf Q/\mathbf Z)_{n,v}\), the quotient \(A\), and the large exactness diagram were reconstructed semantically. The upper row is the Weil-group sequence, the lower row is the Galois/Hasse sequence, the two right-hand \((\mathbf Q/\mathbf Z)_n\) terms are identified, the middle vertical quotient has kernel \(\bigoplus_v(\mathbf Q/\mathbf Z)_{n,v}\), and the two copies of \(A\) are identified above. The inherited TeX changes this topology and was not promoted. In 5.1, orthogonal direct sums and the totally isotropic filtration are preserved exactly. No crop.\par

\subsection*{Authority page 16 / printed folio 314}

Material comparison decisions. The authority continuation says that \(\mathscr V\), not \(\mathscr W^\perp\), is isomorphic to \((\mathscr W\oplus\mathscr V/\mathscr W^\perp)\oplus(\mathscr W^\perp/\mathscr W)\). The Chern-class identity is \(w^2(V)=c^1(W)\bmod 2\), not the inherited \(\operatorname{cl}(W)\). Proposition 5.2, the epsilon identity, the semisimple decomposition, Corollary 5.3, and the start of 5.4 were replayed from the authority. No source defect required repair and no crop was used.\par

\subsection*{Authority page 17 / printed folio 315}

Canonical repairs and formula control. The scan’s grammatical omission “il n’en pas de même” is repaired to “il n’en est pas de même”. The displayed formula belongs to subsection 5.4 and is canonically numbered (5.4.1), replacing the printed (2.3.1); the inherited candidate’s (5.3.1) is also rejected. Standard Weil--Deligne typing requires the associated pair \((\rho',N)\), not the scan’s isolated \((\rho',N')\), because item (b), the covariance relation, the determinant, and the lemma all use \(N\). The opening \(\varepsilon(V)=\varepsilon(V,0)\) is an evaluation at \(s=0\) and is retained, as is the statement that \(t\) is morally \(q^{-s}\); inherited substitutions of an additive character or of \(q^{-1/2}\) at that point were rejected. No crop.\par

\subsection*{Authority page 18 / printed folio 316}

End matter. Bibliography entries 1--3, the received date “Reçu le 20 septembre 1975”, the author’s name, I.H.E.S., street address, postal code, and country were transcribed from the authority. The English record translates only the received-date label; bibliographic titles and the address remain as printed. The large terminal blank remainder, running head, and folio were excluded. No crop.\par

\end{document}
